% Section 13
\section{Appendix E: Regularisation of the Aichelburg--Sexl Singularity}
\label{sec:as_reg}

The derivation of $\Delta A$ in Section~\ref{sec:deltaA} applies to 
a normalizable photon mode but inherits from the strict 
Aichelburg--Sexl construction a distributional stress tensor 
$T_{uu} \propto \delta^{(2)}(\mathbf{x}_\perp)$ and a longitudinal 
shift $\Delta v \propto -\ln r^2$ that diverges on the axis $r \to 0$. 
The electromagnetic construction of Section~\ref{sec:microstates} 
provides a natural regularisation: the photon is a normalizable 
mode of the electromagnetic field in double-null gauge, with a 
smooth Gaussian transverse profile of characteristic width 
$\sigma \sim c/\omega$ set by its wavelength. We carry out the 
regularisation explicitly here.

\paragraph{Photon stress tensor and energy conservation.}
In double-null gauge with $A_+ = A_- = 0$, the only nonzero field 
strengths at the shock front are $F_{ui} = \partial_u A_i$. Since 
$g_{uu} = 0$ in double-null coordinates, the stress tensor 
component sourcing the shift is
%
\begin{equation}
    T_{uu}(u, \mathbf{x}_\perp) 
    = \sum_{i=1,2}(\partial_u A_i)^2 
    = \hbar\omega\, f(u)\, g_\sigma(r),
    \label{eq:Tuu_reg}
\end{equation}
%
where $f(u)$ is a normalised longitudinal profile 
($\int_{-\infty}^{\infty} f\,du = 1$) encoding the photon 
wavepacket, and
%
\begin{equation}
    g_\sigma(r) 
    = \frac{1}{\pi\sigma^2}\,e^{-r^2/\sigma^2}, 
    \qquad \sigma \sim \frac{c}{\omega},
    \label{eq:gaussian_profile}
\end{equation}
%
is a Gaussian transverse profile normalised so that 
$\int g_\sigma\, d^2x_\perp = 1$. 
The total null energy follows immediately from energy conservation:
%
\begin{equation}
    \int T_{uu}\,du\,d^2x_\perp 
    = \hbar\omega \int f\,du \int g_\sigma\,d^2x_\perp 
    = \hbar\omega,
    \label{eq:null_energy_reg}
\end{equation}
%
independently of the profile shape.

\paragraph{Area increment.}
Substituting equation~\eqref{eq:null_energy_reg} into the 
Raychaudhuri result~\eqref{eq:deltaA}:
%
\begin{equation}
    \Delta A 
    = 8\pi G \int T_{uu}\,du\,d^2x_\perp 
    = 8\pi G\hbar\omega.
    \label{eq:deltaA_reg}
\end{equation}
%
The area increment is finite, independent of $\sigma$, and 
determined entirely by the photon energy $\hbar\omega$. The 
regularisation does not alter the area result; it only 
regularises the spatial distribution of the shift.

\paragraph{The regulated longitudinal shift.}
Applying the Einstein equations to the perturbed pp-wave metric 
$ds^2 = -2\,du\,dv + h(\mathbf{x}_\perp)\delta(u)\,du^2 
+ d^2\mathbf{x}_\perp$ at the shock front, the Ricci component 
$R_{uu} = 8\pi G T_{uu}$ gives the two-dimensional Poisson 
equation for the shift $\Delta v = h$:
%
\begin{equation}
    \nabla^2_\perp\,\Delta v 
    = -16\pi G\hbar\omega\; g_\sigma(r)
    = -\frac{16G\hbar\omega}{\sigma^2}\,e^{-r^2/\sigma^2}.
    \label{eq:poisson_reg}
\end{equation}
%
The total source integral 
$\int \nabla^2_\perp\,\Delta v\,d^2x_\perp = -16\pi G\hbar\omega$ 
matches the strict AS case, with the delta function replaced by 
a smooth Gaussian source. Applying Gauss's theorem in two 
dimensions to the radially symmetric source:
%
\begin{equation}
    \frac{d\,\Delta v}{dr} 
    = -\frac{8G\hbar\omega}{r}
      \left(1 - e^{-r^2/\sigma^2}\right),
    \label{eq:grad_reg}
\end{equation}
%
which vanishes at $r = 0$ (the shift is smooth on axis) and 
approaches $-8G\hbar\omega/r$ for $r \gg \sigma$, recovering 
the AS form. 

Integrating equation~\eqref{eq:grad_reg} and using the 
identity 
$\int_0^x (1 - e^{-t})/t\,dt = \gamma + \ln x + E_1(x)$,
where $\gamma$ is the Euler--Mascheroni constant 
and $E_1(x) = \int_x^\infty e^{-t}/t\,dt$ is the exponential 
integral, with reference scale $r_0 \gg \sigma$ 
defined by $\Delta v(r_0) = 0$:
%
\begin{equation}
    \boxed{
    \Delta v(r) 
    = -4G\hbar\omega\!\left[\ln\frac{r^2}{r_0^2} 
      + E_1\!\!\left(\frac{r^2}{\sigma^2}\right)\right].
    }
    \label{eq:shift_reg}
\end{equation}

\paragraph{Limiting behaviour.}

\noindent\textit{On axis, $r \to 0$.}\;
Using the small-argument expansion 
$E_1(x) \to -\gamma - \ln x$ as $x\to 0^+$, the bracket 
in~\eqref{eq:shift_reg} has the finite limit
%
\begin{align}
    \ln\frac{r^2}{r_0^2} + E_1\!\!\left(\frac{r^2}{\sigma^2}\right)
    &\;\xrightarrow{r\,\to\,0}\;
    \ln\frac{r^2}{r_0^2} - \gamma - \ln\frac{r^2}{\sigma^2}
    = \ln\frac{\sigma^2}{r_0^2} - \gamma,
    \label{eq:bracket_limit}
\end{align}
%
giving the finite on-axis value
%
\begin{equation}
    \Delta v(0) 
    = 4G\hbar\omega
    \!\left[\ln\frac{r_0^2}{\sigma^2} + \gamma\right].
    \label{eq:shift_center}
\end{equation}
%
The logarithmic divergence of the strict AS shift at $r = 0$ is 
resolved: $\Delta v(0)$ grows as $\ln(r_0/\sigma)$ but is finite 
for any $\sigma > 0$.

\medskip
\noindent\textit{Far from axis, $r \gg \sigma$.}\;
For large argument, $E_1(r^2/\sigma^2) \to 0$ exponentially, 
and~\eqref{eq:shift_reg} reduces to
%
\begin{equation}
    \Delta v(r \gg \sigma) 
    \;\approx\; -4G\hbar\omega\ln\frac{r^2}{r_0^2},
    \label{eq:shift_AS}
\end{equation}
%
recovering the strict AS logarithmic profile.

\paragraph{Recovery of the strict limit.}
As $\sigma \to 0$ at fixed $r > 0$, $E_1(r^2/\sigma^2) \to 0$ 
and equation~\eqref{eq:shift_reg} reduces to the strict 
form~\eqref{eq:shift_AS} pointwise for all $r > 0$, while 
$\Delta v(0) \to +\infty$ logarithmically via 
equation~\eqref{eq:shift_center}, recovering the AS singularity 
on-axis.
